\begin{abstract}
\emph{Evenwel v. Abbott}, 578 U.S. 54 (2016), held unanimously that states \emph{may} use total population as the redistricting base, but explicitly left open whether states \emph{must} use citizen voting-age population (citizen VAP) or may adopt it as an alternative. The question is not merely academic: the gap between total population and citizen VAP exceeds five percentage points in approximately 80 of the 435 congressional districts, and the systematic geographic pattern of noncitizen concentration makes the choice of population base a substantial determinant of partisan and urban-rural representation balance.

This paper examines the legal framework governing noncitizen-population redistricting at both the congressional and state legislative levels, the empirical magnitude of the total-population--citizen-VAP divergence in the high-immigration states of California, Texas, New York, and Florida, and the partisan and geographic consequences of switching from total population to citizen VAP as the redistricting base. We also describe BISECT's \texttt{--population-source cvap} implementation, which uses the Census CVAP special tabulation to run the recursive bisection algorithm under a citizen-VAP balance constraint, enabling systematic comparison of plans under the two population definitions.

The paper's central empirical finding is that switching from total population to citizen VAP for congressional redistricting produces a consistently Republican-advantaging shift: urban districts with large noncitizen populations lose population credit, requiring them to absorb additional territory; suburban and exurban districts with small noncitizen populations gain relative standing. The effect is largest in California (where the citizen-VAP--total-population gap exceeds 8 percentage points in Los Angeles County tracts), Texas (South Texas border districts), New York (New York City boroughs), and Florida (Miami-Dade). We estimate that switching to citizen VAP would shift 15--20 congressional seats nationally from urban-Democratic to suburban-Republican lean over a sustained redistricting cycle.

The legal distinction between congressional and state legislative redistricting is critical: \emph{Article I}'s command that Representatives be ``apportioned among the several States according to their respective Numbers'' has consistently been interpreted to require total-population apportionment at the federal level, insulating congressional redistricting from citizen-VAP mandates. State legislative redistricting, addressed by \emph{Evenwel}, presents a different constitutional question. We analyze the arguments for and against citizen-VAP mandates in both contexts.
\end{abstract}
